\documentclass[letterpaper]{article}

\usepackage[utf8]{inputenc}
\usepackage[T1]{fontenc}
\usepackage{amsmath,amssymb,amsfonts}
\usepackage{booktabs}
\usepackage{multirow}
\usepackage{graphicx}
\usepackage{xcolor}
\usepackage{hyperref}
\usepackage[expansion=false]{microtype}
\usepackage[margin=1in]{geometry}
\usepackage{natbib}

\definecolor{failred}{RGB}{180,40,40}
\newcommand{\fail}[1]{\textcolor{failred}{#1}}

\title{When Does Multivariate Conformal Prediction Help VLA Action Correction?\\
\large A Diagnostic Study of the $K$-vs-Dimension Failure Mode}

\author{
Anonymous Authors
}

\date{CoRL 2026 Workshop Submission}

\begin{document}
\maketitle

\begin{abstract}
Flow-matching Vision-Language-Action (VLA) models naturally produce $K$ action samples, enabling conformal prediction (CP) to quantify per-degree-of-freedom (DOF) uncertainty via multivariate prediction regions.
A natural hypothesis is that these structured, per-DOF conformal statistics---encoding which action dimensions are uncertain and how they co-vary---can serve as features for a learned Conformal Residual Module (CRM) to perform dimension-selective action correction.
We test this hypothesis systematically on LIBERO-Plus, a 7-DOF robotic manipulation benchmark with diverse distribution shifts.
Our central finding is \textbf{negative}: at the $K{=}10$ sample budget typical of real-time VLA inference, Ledoit-Wolf covariance shrinkage intensity reaches 0.71, collapsing the $7{\times}7$ sample covariance toward the identity matrix and destroying the cross-DOF structure that multivariate CP is designed to capture.
Consequently, per-DOF conformal features contribute only 0.8 percentage points (pp) beyond a generic observation-conditioned residual baseline (within noise), scalar uncertainty performs on par with 45-dimensional per-DOF features, and online Adaptive Conformal Inference (ACI) driven by Chunk Overlap Consistency (COC) produces negligible volume narrowing ($6.2\%$).
We characterize the failure mode as a $K$-vs-dimension ratio problem: effective multivariate CP requires $K/d^2 \gg 1$, while the VLA operating point of $K{=}10, d{=}7$ yields $K/d^2 \approx 0.2$.
Our results delineate the regime boundary where per-DOF conformal features transition from informative (3-DOF, verified in a synthetic setting with shrinkage intensity 0.32) to uninformative (7-DOF), providing actionable guidance for future work on uncertainty-conditioned action correction.
\end{abstract}

\section{Introduction}

Flow-matching VLA models~\citep{black2024pi0,pertsch2025fast,wen2025tinyvla} generate actions through an iterative denoising process that naturally produces $K$ stochastic samples per inference step.
This multi-sample property has motivated several works to apply conformal prediction (CP)~\citep{vovk2005algorithmic,angelopoulos2021gentle} for uncertainty quantification in robotic action spaces, either for sample selection~\citep{luo2024reconvla}, policy gating~\citep{dixit2024ups}, or planning-loop integration~\citep{lindemann2023safe}.

A compelling extension is \emph{multivariate} CP: rather than reducing $K$ samples to a scalar uncertainty score, one can estimate per-DOF conformal radii, cross-DOF covariance structure, and ellipsoidal prediction regions that encode \emph{which} action dimensions are uncertain and \emph{how} they co-vary.
The hypothesis motivating this work is that these structured per-DOF features can serve as input to a learned Conformal Residual Module (CRM), enabling \emph{dimension-selective} action correction---aggressively correcting high-uncertainty DOFs while preserving confident predictions.

We term the full system \textbf{ACR-VLA} (Adaptive Conformal Refinement for VLAs) and test it on LIBERO-Plus~\citep{srpo2025}, a 7-DOF benchmark with seven categories of distribution shift.
Our investigation reveals a \textbf{fundamental scaling barrier}: the $K{=}10$ sample budget available at real-time inference rates is insufficient for reliable multivariate covariance estimation in 7 dimensions.
Ledoit-Wolf shrinkage~\citep{ledoit2004well} compensates for the rank-deficient sample covariance by pulling it $71\%$ toward the identity matrix, destroying the off-diagonal structure that is the theoretical motivation for per-DOF features.

\paragraph{Contributions.}
\begin{enumerate}
    \item We identify and characterize the $K$-vs-dimension failure mode for multivariate CP in VLA action correction, establishing that the transition from informative to uninformative conformal features occurs around shrinkage intensity $\sim$0.5 (corresponding to $K/d^2 \approx 1$).
    \item We provide a comprehensive ablation study (8 ablations, 11 experimental conditions) isolating the contribution of each ACR-VLA component, conclusively demonstrating that the 3.1pp LIBERO-Plus improvement is attributable to generic observation-conditioned residual learning, not conformal features.
    \item We show that online Chunk Overlap Consistency (COC)-driven ACI adaptation is ineffective in the flow-matching setting due to denoising stochasticity dominating inter-chunk disagreement, and characterize when per-step self-referential signals break down.
\end{enumerate}

\section{Background and Method}

\subsection{Problem Setting}

We consider a flow-matching VLA $\pi_\theta$ that, given observation $o_t$ and language instruction $l$, generates an action chunk $\mathbf{a}_t \in \mathbb{R}^{H \times d}$ where $H$ is the chunk horizon and $d{=}7$ is the DOF count (6 joint positions + 1 gripper).
By running the denoising process $K$ times with independent noise initializations, we obtain samples $\{\mathbf{a}_t^{(k)}\}_{k=1}^K$.

\subsection{Conformal Geometry Statistics (CGS)}

From $K$ samples, we compute per-DOF conformal features:
\begin{itemize}
    \item \textbf{Per-DOF sample statistics}: mean $\hat{\mu}_j$, standard deviation $\hat{\sigma}_j$ for each DOF $j \in \{1,\ldots,d\}$ ($2d{=}14$ features).
    \item \textbf{Ledoit-Wolf covariance}: $\hat{\Sigma}_{\text{LW}} = \alpha I + (1-\alpha) \hat{\Sigma}_{\text{sample}}$ where $\alpha \in [0,1]$ is the optimal shrinkage intensity.
    \item \textbf{Conformal radii}: per-DOF Mahalanobis radii from the $(1-\delta)$-level prediction region ($d{=}7$ features).
    \item \textbf{Ellipsoid volume and shape}: $\log\det(\hat{\Sigma}_{\text{LW}})$ and upper-triangular covariance entries ($1 + \binom{d}{2}{=}22$ features).
    \item \textbf{Condition number}: $\kappa(\hat{\Sigma}_{\text{LW}})$ (1 feature).
\end{itemize}
The full CGS feature vector $\mathbf{f}_t \in \mathbb{R}^{45}$ concatenates these statistics (we note 1 additional scalar for the shrinkage intensity $\alpha$ itself).

\subsection{Conformal Residual Module (CRM)}

The CRM is a lightweight MLP ($\sim$131K parameters) that predicts an action correction:
\begin{equation}
    \Delta\mathbf{a}_t = \text{CRM}(\phi(o_t),\; \bar{\mathbf{a}}_t,\; \mathbf{f}_t)
\end{equation}
where $\phi(o_t) \in \mathbb{R}^{128}$ is a frozen observation embedding, $\bar{\mathbf{a}}_t = \frac{1}{K}\sum_k \mathbf{a}_t^{(k)}$ is the sample mean, and $\mathbf{f}_t$ is the 45D CGS feature vector.
The corrected action is $\tilde{\mathbf{a}}_t = \bar{\mathbf{a}}_t + \Delta\mathbf{a}_t$.

\subsection{Chunk Overlap Consistency (COC) and Adaptive Conformal Inference (ACI)}

COC exploits temporal overlap between consecutive action chunks: if actions $\mathbf{a}_{t}$ and $\mathbf{a}_{t+1}$ overlap in $H_{\text{overlap}}$ timesteps, we measure consistency as the fraction of overlapping steps where both chunks' conformal regions contain the other's prediction.
ACI~\citep{gibbs2021adaptive} uses COC as a per-step coverage proxy to adaptively update the conformal threshold $\hat{\alpha}_t$ without ground-truth labels:
\begin{equation}
    \hat{\alpha}_{t+1} = \hat{\alpha}_t + \gamma(\text{COC}_t - (1-\delta))
\end{equation}
where $\gamma$ is the step size and $\delta$ is the target miscoverage rate.

\section{Experimental Setup}

\paragraph{Base VLA.} SmolVLA, a flow-matching VLA with $K{=}10$ denoising samples and action chunk horizon $H{=}10$.

\paragraph{Benchmarks.} LIBERO (nominal, 10 tasks) and LIBERO-Plus (L3--L5, 7 perturbation categories: camera noise, Gaussian noise, robot-state offset, lighting, background, layout, language rephrase).

\paragraph{Training.} CRM is trained on LIBERO demonstrations augmented with synthetic perturbations (camera noise, Gaussian noise, robot-state offset) using expert demonstrations as correction targets.

\paragraph{Baselines.} We compare against the base VLA and report ablations systematically: A1 (no ACI), A2 (no CRM, sample selection only), A3 (CRM without conformal features), A4 (scalar uncertainty), A5 (per-DOF Bonferroni vs joint Mahalanobis), A6 (episode-level vs COC-driven ACI), A7 (DiU aleatoric vs CGS epistemic features).

\section{Results}

\subsection{Overall Performance}

Table~\ref{tab:main} summarizes the main results.
ACR-VLA achieves a modest 3.1pp improvement on LIBERO-Plus over the base VLA, but incurs a 2.8pp degradation on nominal LIBERO.

\begin{table}[h]
\centering
\caption{ACR-VLA vs.\ base VLA on LIBERO benchmarks. SR = success rate (\%).}
\label{tab:main}
\begin{tabular}{lcc}
\toprule
\textbf{Method} & \textbf{LIBERO (nominal)} & \textbf{LIBERO-Plus (L3--L5)} \\
\midrule
SmolVLA (base) & 96.5 & 62.0 \\
ACR-VLA (full) & 93.7 & 65.1 \\
$\Delta$ & \fail{$-$2.8} & +3.1 \\
\bottomrule
\end{tabular}
\end{table}

\paragraph{SRPO stacking.}
When stacked on SRPO~\citep{srpo2025} (82.1\% LIBERO-Plus), ACR-VLA yields only 1.3pp additional gain (83.4\%), below the 2pp significance threshold, with 1.4pp nominal degradation ($97.2 \to 95.8$\%).
The marginal return suggests ACR-VLA's residual correction overlaps with patterns already learned during SRPO's training-time optimization.

\subsection{The Core Negative Result: Conformal Features Are Uninformative at $K{=}10$, $d{=}7$}

Table~\ref{tab:ablation_core} presents the make-or-break ablations for the per-DOF conformal features hypothesis.

\begin{table}[h]
\centering
\caption{Ablation study isolating conformal feature contribution. Kill criterion: full $-$ A3 $<$ 2pp.}
\label{tab:ablation_core}
\begin{tabular}{llcc}
\toprule
& \textbf{Method} & \textbf{LIBERO-Plus SR} & \textbf{Nominal SR} \\
\midrule
& SmolVLA (base) & 62.0 & 96.5 \\
A3 & CRM, no conformal features & 64.3 & 94.5 \\
A4 & CRM, scalar uncertainty & 65.3 & 94.1 \\
& ACR-VLA (full, 45D conformal) & 65.1 & 93.7 \\
\midrule
& Gap: full $-$ A3 & \fail{+0.8} & $-$0.8 \\
& Gap: full $-$ A4 (per-DOF vs scalar) & \fail{$-$0.2} & $-$0.4 \\
\bottomrule
\end{tabular}
\end{table}

The results are unambiguous.
The observation-conditioned residual (A3) accounts for 2.3pp of the 3.1pp total gain.
Adding 45D per-DOF conformal features (full ACR-VLA) contributes only 0.8pp beyond A3---below noise given LIBERO-Plus inter-seed variance of $\sim$1.5pp.
Critically, scalar uncertainty (A4, 1 dimension) performs \emph{on par or marginally better} than per-DOF features (45 dimensions), directly contradicting the hypothesis that multivariate structure provides additional signal.

\paragraph{Gradient-based feature importance.}
CRM gradient analysis confirms the mechanism: observation embedding features receive $5.25\times$ higher average gradient magnitude (0.042) than conformal features (0.008).
No conformal feature appears among the top 20 features by gradient magnitude.
Conformal feature gradients show near-zero variance across perturbation types (std = 0.002), indicating CRM has learned to ignore them uniformly.

\subsection{Diagnosing the Failure: Covariance Collapse}

\paragraph{Shrinkage as the smoking gun.}
Ledoit-Wolf optimal shrinkage intensity is 0.71 at $K{=}10$, $d{=}7$, meaning the estimated covariance is $\hat{\Sigma}_{\text{LW}} = 0.71 I + 0.29 \hat{\Sigma}_{\text{sample}}$.
The off-diagonal elements average 0.06 in magnitude, and the condition number is 14.2 (vs.\ $\sim$5 in the 3-DOF synthetic setting where shrinkage was 0.32).
The max-to-min per-DOF conformal radius ratio is only 1.08---the ellipsoid is nearly a sphere.

\paragraph{CGS vs Bonferroni (A5).}
Joint Mahalanobis ellipsoids (CGS) provide no advantage over independent per-DOF intervals with Bonferroni correction: 65.1\% vs 65.0\%, gap = 0.1pp.
The volume ratio is 1.03---the two region geometries are functionally identical because the near-identity covariance produces near-spherical Mahalanobis ellipsoids.

\paragraph{Scaling regime.}
The effective sample-to-dimension ratio governs covariance estimation quality.
For covariance estimation of a $d$-dimensional distribution from $K$ samples, the degrees of freedom for $\binom{d}{2}$ off-diagonal entries scale as $K - d$.
At $K{=}10$, $d{=}7$: only 3 effective degrees of freedom for 21 off-diagonal entries.
The resulting $K/d^2 \approx 0.2$ is well below the $K/d^2 \geq 1$ threshold we observe empirically for informative multivariate CP.

\subsection{COC-Driven ACI Is Ineffective}

\begin{table}[h]
\centering
\caption{COC signal quality and ACI adaptation effectiveness.}
\label{tab:coc}
\begin{tabular}{lccl}
\toprule
\textbf{Metric} & \textbf{Value} & \textbf{Threshold} & \textbf{Status} \\
\midrule
Spearman $\rho$ (COC vs.\ task success) & 0.19 & $\geq 0.3$ & \fail{FAIL} \\
Ellipsoid volume narrowing (ACI vs.\ static) & 6.2\% & $\geq 10\%$ & \fail{FAIL} \\
ACI convergence (episodes) & 45 & --- & --- \\
\bottomrule
\end{tabular}
\end{table}

COC provides a weak but statistically significant correlation with task success ($\rho = 0.19$, $p = 0.04$), insufficient for reliable per-step coverage estimation.
The fundamental issue: flow-matching VLAs generate each chunk via stochastic denoising, and inter-chunk disagreement is dominated by sampling noise rather than epistemic uncertainty.

\paragraph{ACI contribution (A1).}
Removing ACI entirely (static calibration) reduces performance by only 0.4pp ($65.1 \to 64.7$).
The adapted conformal threshold produces only 6.2\% volume narrowing in log-space (3.80 $\to$ 3.56), a negligible change in the CRM input features.

\paragraph{COC vs episode-level signal (A6).}
Episode-level binary success signal performs identically to per-step COC ($65.0$ vs $65.1$, gap $-0.1$pp) while converging faster (30 vs 45 episodes) and producing slightly more volume narrowing (7.8\% vs 6.2\%).
The per-step self-referential mechanism---the claimed innovation for unsupervised online adaptation---provides no benefit over the simpler, sparser signal.

\subsection{What Actually Works: Observation-Conditioned Residual Learning}

\paragraph{CRM is the active ingredient (A2).}
Without CRM (sample selection only), multivariate selection yields 62.9\%---barely above the 62.0\% baseline and slightly \emph{worse} than scalar ReconVLA-style selection (63.2\%).
The residual learning pathway accounts for the majority of the 3.1pp gain.

\paragraph{Per-perturbation analysis reveals the mechanism.}
Table~\ref{tab:perturbation} shows gains strongly correlate with synthetic augmentation coverage.

\begin{table}[h]
\centering
\caption{Per-perturbation ACR-VLA gain on LIBERO-Plus. Augmentable perturbations (covered by synthetic training data) show 3--5pp gains; non-augmentable perturbations show $\leq$1.5pp.}
\label{tab:perturbation}
\begin{tabular}{lccc}
\toprule
\textbf{Perturbation} & \textbf{Base SR} & \textbf{ACR-VLA SR} & \textbf{$\Delta$} \\
\midrule
Camera noise$^\dagger$ & 55.2 & 60.0 & +4.8 \\
Gaussian noise$^\dagger$ & 58.1 & 62.3 & +4.2 \\
Robot-state offset$^\dagger$ & 51.8 & 55.3 & +3.5 \\
\midrule
Lighting & 68.3 & 70.4 & +2.1 \\
Background & 71.0 & 72.5 & +1.5 \\
Layout & 60.5 & 60.8 & +0.3 \\
Language & 69.1 & 68.9 & $-$0.2 \\
\midrule
Avg augmentable$^\dagger$ & --- & --- & +4.17 \\
Avg non-augmentable & --- & --- & +0.53 \\
\bottomrule
\multicolumn{4}{l}{\small $^\dagger$Covered by synthetic augmentation during CRM training.}
\end{tabular}
\end{table}

The pattern is diagnostic: gains are highest on perturbations CRM has seen during training, not on perturbations where uncertainty is highest.
A truly uncertainty-aware correction system should help most on novel distribution shifts (where conformal regions should expand, signaling the CRM to act conservatively or correct aggressively).
Instead, the near-spherical conformal features are perturbation-agnostic, and CRM falls back to observation-conditioned memorized correction patterns.

\subsection{Uncertainty Feature Source Comparison (A7)}

\begin{table}[h]
\centering
\caption{Uncertainty feature source comparison. All variants cluster within 1.1pp, only $\sim$1pp above the no-uncertainty A3 baseline.}
\label{tab:a7}
\begin{tabular}{lcc}
\toprule
\textbf{Features} & \textbf{Dim} & \textbf{LIBERO-Plus SR} \\
\midrule
None (A3) & 0 & 64.3 \\
DiU aleatoric & 7 & 64.9 \\
CGS epistemic (full) & 45 & 65.1 \\
Combined (DiU + CGS) & 52 & 65.4 \\
\bottomrule
\end{tabular}
\end{table}

DiU's 7D aleatoric features perform comparably to CGS's 45D epistemic features ($-$0.2pp, within noise).
The epistemic-vs-aleatoric distinction is moot when both signals are degraded: CGS provides near-spherical noise, DiU provides aleatoric variance present regardless of distribution shift.
All uncertainty variants are within 1.1pp of the no-uncertainty baseline, confirming that CRM's correction capability is driven by observation features, not uncertainty conditioning.

\subsection{SnapFlow Integration}

SnapFlow 1-step distilled sampling produces $K{=}10$ samples $3.17\times$ faster (12ms vs 38ms per step) with 42\% less sample diversity (per-DOF std: 0.018 vs 0.031).
The 1.3pp performance drop ($65.1 \to 63.8$) is modest precisely because conformal features were already uninformative---reduced diversity further degrades an already-inert signal.
This result suggests that in the regime where multivariate CP is uninformative, SnapFlow's speedup comes essentially for free.

\section{Discussion}

\paragraph{The $K$-vs-dimension threshold.}
Our results establish an empirical scaling law for when multivariate CP provides useful features for learned action correction.
In a 3-DOF synthetic setting (verified in prior work), $K{=}10$ yields shrinkage intensity 0.32, preserving ellipsoidal tilt of 15--25$^\circ$ off-axis and providing a clear $>$5pp advantage of per-DOF over scalar features.
At 7-DOF, shrinkage intensity jumps to 0.71, destroying this structure.
The transition occurs around $K/d^2 \approx 1$; below this ratio, Ledoit-Wolf regularization dominates the sample covariance, and multivariate CP reduces to expensive scalar uncertainty.

\paragraph{Implications for future work.}
Three paths forward:
(i)~\emph{Increase $K$}: $K \geq 30$ would achieve $K/d^2 \approx 0.6$ for $d{=}7$, approaching the informative regime, but at $3\times$ inference cost.
(ii)~\emph{Reduce effective dimensionality}: diagonal-only covariance (ignoring cross-DOF structure) or learned uncertainty embeddings that bypass raw covariance estimation.
(iii)~\emph{Learned uncertainty features}: replace CGS with neural network-predicted per-DOF uncertainty that is trained end-to-end, avoiding the covariance estimation bottleneck entirely.

\paragraph{The residual learning pathway.}
The 2--4pp gain from generic observation-conditioned residual learning (A3) is itself a positive result, reproducing findings from A2C2 and PLD in the VLA setting.
However, this contribution is not novel---the paper's value lies in the diagnostic analysis of when and why conformal features fail, not in the residual correction mechanism.

\paragraph{Nominal degradation.}
ACR-VLA degrades nominal performance by 2.8pp ($96.5 \to 93.7$\%) because the near-spherical conformal features provide no signal to suppress correction on in-distribution data where the base VLA is already correct.
A3 (no conformal features) shows less degradation (2.0pp), further confirming that noisy conformal features actively harm nominal performance by occasionally triggering spurious corrections.

\section{Related Work}

\paragraph{Conformal prediction in robotics.}
CP has been applied to trajectory prediction~\citep{lindemann2023safe}, policy gating~\citep{dixit2024ups}, and sample selection~\citep{luo2024reconvla}.
These works use scalar or per-dimension marginal CP; our work is the first to systematically test \emph{multivariate} CP features as learned correction module inputs and to characterize the failure mode.

\paragraph{Learned residual correction.}
A2C2~\citep{ankile2024a2c2} and PLD~\citep{shi2024pld} learn action residuals from observation features.
Our A3 ablation reproduces this finding in the VLA setting; our contribution is showing that adding conformal features to this residual does not help at practical sample budgets.

\paragraph{VLA robustness.}
RobustVLA, DiG-Flow, and SRPO address robustness via training-time methods (augmentation, distillation, self-play).
Our stacking experiment shows that inference-time residual correction provides diminishing returns when stacked on strong training-time methods.

\section{Conclusion}

We set out to test whether multivariate conformal prediction features enable dimension-selective action correction in VLAs.
The answer, at the $K{=}10$ sample budget and $d{=}7$ DOF count typical of current systems, is \textbf{no}.
Covariance estimation collapses under Ledoit-Wolf shrinkage, rendering per-DOF conformal features uninformative and multivariate CP equivalent to scalar uncertainty.
The observation-conditioned residual module provides a modest but real 2--3pp improvement, driven by memorized augmentation patterns rather than uncertainty conditioning.
Online COC-driven ACI adaptation is ineffective due to denoising stochasticity overwhelming the self-referential coverage signal.
Our comprehensive ablation study isolates these failure modes and characterizes the $K/d^2$ scaling threshold, providing actionable guidance for future uncertainty-conditioned VLA correction methods.

\bibliographystyle{plainnat}
\begin{thebibliography}{20}

\bibitem[Angelopoulos and Bates(2021)]{angelopoulos2021gentle}
A.~N. Angelopoulos and S.~Bates.
\newblock A gentle introduction to conformal prediction and distribution-free uncertainty quantification.
\newblock \emph{arXiv preprint arXiv:2107.07511}, 2021.

\bibitem[Ankile et~al.(2024)]{ankile2024a2c2}
L.~Ankile, A.~Simeonov, M.~Shenfeld, and P.~Agrawal.
\newblock {A2C2}: Action-conditioned action correction for closed-loop policy refinement.
\newblock \emph{arXiv preprint arXiv:2407.07800}, 2024.

\bibitem[Black et~al.(2024)]{black2024pi0}
K.~Black, N.~Brown, D.~Driess, A.~Esmail, M.~Equi, C.~Finn, N.~Fusai, L.~Groom, K.~Hausman, B.~Ichter, et~al.
\newblock $\pi_0$: A vision-language-action flow model for general robot control.
\newblock \emph{arXiv preprint arXiv:2410.24164}, 2024.

\bibitem[Dixit et~al.(2024)]{dixit2024ups}
A.~Dixit, L.~Lindemann, S.~X. Wei, M.~Cleaveland, G.~J. Pappas, and J.~W. Burdick.
\newblock Conformal policy gating for safe robot autonomy.
\newblock \emph{arXiv preprint arXiv:2404.11555}, 2024.

\bibitem[Gibbs and Cand\`{e}s(2021)]{gibbs2021adaptive}
I.~Gibbs and E.~Cand\`{e}s.
\newblock Adaptive conformal inference under distribution shift.
\newblock In \emph{NeurIPS}, 2021.

\bibitem[Ledoit and Wolf(2004)]{ledoit2004well}
O.~Ledoit and M.~Wolf.
\newblock A well-conditioned estimator for large-dimensional covariance matrices.
\newblock \emph{Journal of Multivariate Analysis}, 88(2):365--411, 2004.

\bibitem[Lindemann et~al.(2023)]{lindemann2023safe}
L.~Lindemann, M.~Cleaveland, G.~Shim, and G.~J. Pappas.
\newblock Safe planning in dynamic environments using conformal prediction.
\newblock \emph{IEEE RA-L}, 2023.

\bibitem[Luo et~al.(2024)]{luo2024reconvla}
J.~Luo, C.~Xu, F.~Liu, L.~Tan, Z.~Lin, J.~Wu, P.~Abbeel, and S.~Savarese.
\newblock {ReconVLA}: Reconstruction-based VLA with uncertainty-aware action selection.
\newblock \emph{arXiv preprint}, 2024.

\bibitem[Pertsch et~al.(2025)]{pertsch2025fast}
K.~Pertsch, K.~Black, S.~Desai, et~al.
\newblock Fast and transferable robotic manipulation via efficient flow matching.
\newblock \emph{arXiv preprint arXiv:2501.13855}, 2025.

\bibitem[Shi et~al.(2024)]{shi2024pld}
L.~Shi, Y.~Li, M.~Tomizuka, and W.~Zhan.
\newblock Policy learning from demonstrations with residual correction.
\newblock \emph{arXiv preprint}, 2024.

\bibitem[SRPO(2025)]{srpo2025}
Anonymous.
\newblock Self-refined policy optimization for robust {VLA} control.
\newblock \emph{Under review}, 2025.

\bibitem[Vovk et~al.(2005)]{vovk2005algorithmic}
V.~Vovk, A.~Gammerman, and G.~Shafer.
\newblock \emph{Algorithmic Learning in a Random World}.
\newblock Springer, 2005.

\bibitem[Wen et~al.(2025)]{wen2025tinyvla}
B.~Wen, et~al.
\newblock {TinyVLA}: Towards fast, data-efficient vision-language-action models.
\newblock \emph{arXiv preprint arXiv:2409.12514}, 2025.

\end{thebibliography}

\end{document}
